% Copyright 2026 Sreeram Anil
% SPDX-License-Identifier: CC-BY-4.0
% =============================================================================
%  Missions, sensors and fault scenarios — the benchmark data
% =============================================================================
\chapter{Missions, sensors and fault scenarios}
\label{ch:datasets}

% Figures produced by scripts/make_figures.py from the recorded benchmark run.
% The guard keeps the document compilable before the figures have been regenerated.
\providecommand{\estkitfig}[3]{%
  \begin{figure}[H]\centering
  \IfFileExists{figures/#1}%
    {\includegraphics[width=0.72\textwidth]{figures/#1}}%
    {\fbox{\parbox{0.86\textwidth}{\vspace{8mm}\centering\itshape This figure is regenerated by
      \texttt{scripts/make\_figures.py} from the benchmark run recorded in
      \texttt{results/final}.\vspace{8mm}}}}%
  \caption{#2}\label{#3}\end{figure}}

A benchmark is only as informative as the data it runs on. Seventy estimators evaluated on a single
well-behaved discharge curve would rank almost identically and teach nothing; the same estimators
evaluated on data containing the specific pathologies each method was invented to defeat produce a
ranking that is a design guide. This chapter specifies every dataset used in this report: four
mission profiles, the sensor chain that observes them, twelve cell-level fault and operating
scenarios, a twelve-cell pack dataset, a nine-axis inertial flight dataset, and a 120-flight ageing
protocol. All of it is generated from the plants of Chapters~\ref{ch:ecm} and~\ref{ch:spm} by
\code{battery/datasets.hpp}, \code{pack/pack\_dataset.hpp} and \code{attitude/imu\_dataset.hpp},
deterministically from an integer seed.

Everything in this chapter is \emph{synthetic}. That is the central methodological trade of this
work: we give up the realism of measured cells and buy, in exchange, an exact ground truth, faults
switched on one at a time with known magnitude, and a dataset any reader can regenerate
bit-for-bit. Section~\ref{sec:ds-truth} states precisely which quantities are hidden from the
estimators, because a benchmark in which the estimator can see the answer is worthless.

\section{Design requirements for the mission profiles}

An electric vertical take-off and landing (eVTOL) aircraft loads its battery in a way no road
vehicle does. Yang et~al.~\cite{yang2021evtol} and Fredericks et~al.~\cite{fredericks2018vtol}
describe the power signature of the lift-plus-cruise and tilt-rotor configurations: a short,
extremely high-power hover for take-off, a brief transition during which the wing unloads the
rotors, a moderate-power climb, a long low-power cruise, then the whole sequence in reverse, ending
with a second hover as demanding as the first but flown from a battery already 70--80\,\%
discharged. Three consequences drive the design of the profiles used here.

\emph{The worst-case power draw occurs at the worst-case state of charge}: the landing hover is the
design point for both the pack and the estimator, because the reserve computation --- how many
seconds of hover remain --- depends directly on the state-of-charge (SOC) estimate, and
DO-311A~\cite{rtca2017do311a} requires that indication to be substantiated, which in practice means
an error bound rather than a number. \emph{The C-rates are three to five times automotive}: hover at
3--5\,C (one ``C'' being the current that discharges the nominal capacity in one hour) means the
overpotential the model must subtract from the terminal voltage is several hundred millivolts, an
order of magnitude larger than the open-circuit-voltage (OCV) change over a percent of SOC, so model
error translates into SOC error with a large gain during exactly the phases that matter. And
\emph{the excitation is non-stationary}: hover, transition, climb and cruise excite the cell at
completely different frequencies, so a filter tuned on cruise data can be badly mistuned in hover.
The profiles keep those phases adjacent so that the benchmark sees the transitions.

The cell is a \SI{5}{\ampere\hour} NMC811/graphite 21700-class high-power cell
(Chapter~\ref{ch:ecm}); all C-rates below are referred to its \SI{5}{\ampere\hour} nominal
capacity, and the sign convention throughout the library is \emph{discharge positive}. The sample
rate is \SI{10}{\hertz} ($\dt=\SI{0.1}{\second}$), a typical BMS measurement cadence.

\section{The four mission profiles}
\label{sec:ds-profiles}

A mission profile is a list of segments, each with a duration, a C-rate, a gust fraction and a
label (\code{struct Segment}). The demanded current in segment $s$ at sample $k$ is
\begin{equation}
i_k \;=\; C_s\,Q_{\text{nom}}\,\bigl(1 + \gamma_s\, g_k\bigr)\cdot \pi_k ,
\label{eq:ds-current}
\end{equation}
where $C_s$ is the segment C-rate, $Q_{\text{nom}}=\SI{5}{\ampere\hour}$ the nominal capacity,
$\gamma_s$ the gust fraction, $g_k$ the gust process of \S\ref{sec:ds-gust}, and $\pi_k\in\{1,1.3\}$
the manoeuvre-pulse multiplier of \S\ref{sec:ds-gust}.

\subsection{eVTOL --- the reference mission}

\begin{table}[H]\centering\small
\caption{The \code{evtol} reference mission. C-rates refer to the \SI{5}{\ampere\hour} nominal
capacity; ``\si{\ampere\hour}'' is the nominal (gust-free, pulse-free) charge removed in the
segment. Total \SI{2010}{\second} = \SI{33.5}{\minute}, \SI{3.746}{\ampere\hour} nominal
(\SI{3.763}{\ampere\hour} realised with gusts and manoeuvre pulses at seed~1), a depth of
discharge of \SI{74.9}{\percent} of nominal capacity; the mission starts at $\soc_0=0.95$ and ends
at $\soc\approx0.196$.}
\label{tab:ds-evtol}
\begin{tabular}{llrrrr}\toprule
\# & segment & duration [\si{\second}] & C-rate & current [\si{\ampere}] & charge [\si{\ampere\hour}] \\ \midrule
1  & \code{preflight}     & 120 & 0.20 & 1.0  & 0.033 \\
2  & \code{taxi}          &  60 & 0.50 & 2.5  & 0.042 \\
3  & \code{hover\_takeoff}&  75 & 4.50 & 22.5 & 0.469 \\
4  & \code{transition}    &  45 & 3.00 & 15.0 & 0.188 \\
5  & \code{climb}         & 240 & 2.00 & 10.0 & 0.667 \\
6  & \code{cruise}        & 900 & 1.10 & 5.5  & 1.375 \\
7  & \code{descent}       & 300 & 0.60 & 3.0  & 0.250 \\
8  & \code{transition}    &  45 & 3.00 & 15.0 & 0.188 \\
9  & \code{hover\_landing}&  75 & 4.50 & 22.5 & 0.469 \\
10 & \code{taxi}          &  60 & 0.50 & 2.5  & 0.042 \\
11 & \code{postflight}    &  90 & 0.20 & 1.0  & 0.025 \\ \midrule
   & \textbf{total}       & \textbf{2010} & 1.34 (mean) & 6.71 (mean) & \textbf{3.746} \\
\bottomrule\end{tabular}
\end{table}

The gust fractions are \num{0.05} on the ground, \num{0.08} in hover (the rotors are rate-limited
and the control loop rejects part of the disturbance), \numrange{0.10}{0.12} in transition, climb and
cruise, and \num{0.15} in descent, where the aircraft is lightly loaded and a gust has the largest
relative effect on demanded power. At seed~1 the peak current is \SI{32.0}{\ampere} (6.4\,C) during
the take-off hover including a manoeuvre pulse, the peak cell temperature \SI{36.9}{\celsius} from a
\SI{25}{\celsius} ambient, and the terminal voltage spans \SIrange{3.04}{4.11}{\volt}.

\estkitfig{dataset_evtol.png}{The \code{evtol} reference mission on the ECM plant with nominal
sensors (seed~1): demanded current, true terminal voltage, true SOC and cell temperature. The two
hover segments are the \SI{22.5}{\ampere} spikes at \SI{3}{\minute} and \SI{29}{\minute}; note that
the landing hover drives the terminal voltage below \SI{3.1}{\volt} although the cell still holds
20\,\% SOC --- the \SI{0.27}{\volt} of overpotential at that current is what a voltage-only
estimator has to model away.}{fig:ds-evtol}

\subsection{Fixed-wing --- the long-cruise contrast}

\code{fixed\_wing} represents a conventional take-off and landing electric aircraft: no hover, a
longer cruise at a lower C-rate, a go-around. It is the profile on which Coulomb-counting drift has
the longest time to accumulate and on which OCV-based methods have the least excitation.

\begin{table}[H]\centering\small
\caption{The \code{fixed\_wing} mission. Total \SI{2880}{\second} = \SI{48}{\minute},
\SI{3.942}{\ampere\hour} nominal (\SI{3.959}{\ampere\hour} realised at seed~1), depth of discharge
\SI{78.8}{\percent}, end SOC $\approx 0.157$. Peak current \SI{19.9}{\ampere}, peak cell
temperature \SI{31.9}{\celsius}.}
\label{tab:ds-fixedwing}
\begin{tabular}{llrrrrr}\toprule
\# & segment & duration [\si{\second}] & C-rate & current [\si{\ampere}] & charge [\si{\ampere\hour}] & gust $\gamma_s$ \\ \midrule
1  & \code{preflight} & 120 & 0.20 & 1.0  & 0.033 & 0.05 \\
2  & \code{taxi}      &  90 & 0.60 & 3.0  & 0.075 & 0.10 \\
3  & \code{takeoff}   &  90 & 3.00 & 15.0 & 0.375 & 0.08 \\
4  & \code{climb}     & 420 & 2.00 & 10.0 & 1.167 & 0.10 \\
5  & \code{cruise}    &1500 & 0.80 & 4.0  & 1.667 & 0.12 \\
6  & \code{descent}   & 360 & 0.30 & 1.5  & 0.150 & 0.15 \\
7  & \code{go\_around}&  60 & 2.50 & 12.5 & 0.208 & 0.10 \\
8  & \code{approach}  & 120 & 1.20 & 6.0  & 0.200 & 0.10 \\
9  & \code{taxi}      &  60 & 0.60 & 3.0  & 0.050 & 0.10 \\
10 & \code{postflight}&  60 & 0.20 & 1.0  & 0.017 & 0.05 \\ \midrule
   & \textbf{total}   & \textbf{2880} & 0.99 (mean) & 4.93 (mean) & \textbf{3.942} & \\
\bottomrule\end{tabular}
\end{table}

\estkitfig{dataset_fixed_wing.png}{The \code{fixed\_wing} mission: the \SI{25}{\minute} cruise at
\SI{4}{\ampere} is the most benign excitation in the benchmark and therefore the least
observable.}{fig:ds-fixedwing}

\subsection{Hover-hold --- the thermal worst case}

\code{hover\_hold} is a \SI{19}{\minute} sortie dominated by a \SI{10}{\minute} hover at 3.5\,C ---
a search-and-rescue or inspection mission, and the only profile in which the cell self-heats far
enough for the Arrhenius scaling of the impedance (Chapter~\ref{ch:ecm}) to matter within one
flight. The peak cell temperature is \SI{41.3}{\celsius}, \SI{16}{\kelvin} above ambient, reducing
$R_0$ from \SI{12.0}{\milli\ohm} to about \SI{7.9}{\milli\ohm}. A filter that treats the parameters
as constant sees this as a slow, correlated voltage residual and, if it is adaptive, will absorb it
into the wrong state.

\begin{table}[H]\centering\small
\caption{The \code{hover\_hold} mission. Total \SI{1140}{\second} = \SI{19}{\minute},
\SI{3.950}{\ampere\hour} nominal (\SI{3.997}{\ampere\hour} realised at seed~1), depth of discharge
\SI{79.0}{\percent}, end SOC $\approx 0.149$, mean current \SI{12.5}{\ampere} (2.5\,C).}
\label{tab:ds-hover}
\begin{tabular}{llrrrrr}\toprule
\# & segment & duration [\si{\second}] & C-rate & current [\si{\ampere}] & charge [\si{\ampere\hour}] & gust $\gamma_s$\\ \midrule
1 & \code{preflight} &  60 & 0.20 & 1.0  & 0.017 & 0.05 \\
2 & \code{hover}     & 600 & 3.50 & 17.5 & 2.917 & 0.10 \\
3 & \code{cruise}    & 300 & 1.00 & 5.0  & 0.417 & 0.10 \\
4 & \code{hover}     & 120 & 3.50 & 17.5 & 0.583 & 0.10 \\
5 & \code{post}      &  60 & 0.20 & 1.0  & 0.017 & 0.05 \\ \midrule
  & \textbf{total}   & \textbf{1140} & 2.49 (mean) & 12.47 (mean) & \textbf{3.950} & \\
\bottomrule\end{tabular}
\end{table}

\estkitfig{dataset_hover_hold.png}{The \code{hover\_hold} mission: the \SI{10}{\minute} hover
raises the cell to \SI{41}{\celsius} and lowers the ohmic drop by a third over the course of the
segment.}{fig:ds-hover}

\subsection{Ground charge --- a constant-current/constant-voltage loop}

\code{ground\_charge} is not a flight but the recharge that separates two of them. It exists because
it is the only dataset with negative current, so it exercises the coulombic-efficiency and
hysteresis branches; because it contains a long period of current smoothly \emph{decreasing to
zero}, the condition under which OCV re-anchoring becomes possible; and because it is where a
state-of-health estimator gets its cleanest capacity information.

The profile is a single \SI{3600}{\second} segment with $C_s=-1.00$ (a \SI{5}{\ampere} charge), but
the current is not taken from \eqref{eq:ds-current} directly: the simulator closes a voltage loop
around the plant. If the terminal voltage that the demanded current would produce exceeds
\SI{4.20}{\volt}, a 20-step bisection finds the charge current $i^\star$ for which
$v(i^\star)=\SI{4.20}{\volt}$ and uses that instead; the current is set to zero once its magnitude
falls below \SI{0.25}{\ampere} ($C/20$) or the cell reaches $\soc=0.999$. This is the textbook
CC--CV protocol~\cite{plett2015bms1} implemented plant-in-the-loop rather than as a prescribed
waveform, so it adapts automatically to an aged or cold cell. Starting from $\soc_0=0.20$, the
constant-current phase lasts \SI{2315}{\second} and ends at $\soc=0.841$; the taper then runs
\SI{1096}{\second} until the current drops below $C/20$ at $t=\SI{3411}{\second}$ with
$\soc=0.999$. In total \SI{4.008}{\ampere\hour} are returned and the cell warms only to
\SI{26.7}{\celsius}.

\estkitfig{dataset_ground_charge.png}{The \code{ground\_charge} CC--CV loop: constant
\SI{5}{\ampere} until \SI{4.20}{\volt} at \SI{38.6}{\minute}, then a voltage-limited taper. The
final \SI{3}{\minute} at negligible current is the only genuine rest in the entire cell-level
benchmark.}{fig:ds-charge}

\subsection{Gusts and manoeuvre pulses}
\label{sec:ds-gust}

Real flight-load spectra are broadband but strongly correlated: turbulence models used for flight
qualification (the Dryden and von K\'arm\'an spectra of MIL-HDBK-1797~\cite{milhdbk1797}) are
low-pass with a correlation length set by the turbulence scale divided by the airspeed, giving
correlation times of order seconds for a small aircraft. Reproducing a full Dryden filter would add
parameters without adding insight, so the benchmark uses its discrete first-order approximation: a
stationary Gauss--Markov (AR(1)) process
\begin{equation}
g_k \;=\; a\,g_{k-1} \;+\; \sqrt{1-a^2}\;\varepsilon_k, \qquad
a \;=\; e^{-\dt/\tau_g}, \qquad \tau_g = \SI{2}{\second}, \qquad
\varepsilon_k \sim \N{0}{1},
\label{eq:ds-ar1}
\end{equation}
so that $\E{g_k}=0$, $\Var{g_k}=1$ and $\E{g_k g_{k-m}} = a^{m}$. With $\dt=\SI{0.1}{\second}$,
$a = e^{-0.05} = 0.9512$. The scaling $\sqrt{1-a^2}$ is what makes the \emph{stationary} variance
exactly unity, so $\gamma_s$ in \eqref{eq:ds-current} is directly the relative standard deviation
of the load in that segment: a \num{0.12} gust fraction in cruise means a \SI{12}{\percent} r.m.s.
load fluctuation with a \SI{2}{\second} correlation time. The process is continued across segment
boundaries (it is not reset), so the transition into hover inherits the gust state of the taxi.

On top of the gusts, a \emph{manoeuvre pulse} multiplies the current by \num{1.3} for
\SI{3}{\second}. Pulses are scheduled deterministically at \SI{60}{\second} into a segment and
every \SI{120}{\second} thereafter, and only in segments whose gust fraction exceeds \num{0.07}
(that is, in flight, not on the ground). In the eVTOL mission this yields thirteen pulses --- one
in each hover, two in the climb, seven in the cruise and two in the descent --- adding
\SI{0.027}{\ampere\hour} (\SI{0.55}{\percent}) to the mission total. The pulses are what makes the
load non-Gaussian: they are the step-like, non-zero-mean excursions against which the
innovation-based adaptive filters of Part~IV and the robust filters of Part~V must not overreact.

Both stochastic streams are drawn from a dedicated \code{Rng} instance seeded with
$7919\,s + 17$ for seed $s$, which is independent of the sensor-noise stream (seeded
$104729\,s+3$). Two consequences are worth stating: the \emph{load} realisation is identical for
every estimator and every scenario at a given seed, so differences in the results are never caused
by a different flight; and changing the sensor noise never changes the flight.

\section{The sensor chain}
\label{sec:ds-sensors}

The estimator never sees the plant. It sees three measured signals, produced by the model in
\code{simulate\_mission}:
\begin{align}
v^{\text{m}}_k &= \mathcal{Q}_{v}\!\Bigl(v_k + \sigma_v\,\nu^v_k + o_k + d_k\Bigr),
 & \nu^v_k &\sim \N{0}{1}, \label{eq:ds-vmeas}\\
i^{\text{m}}_k &= \mathcal{Q}_{i}\!\Bigl(\kappa\, i_k + b_i + \sigma_i\,\nu^i_k\Bigr),
 & \nu^i_k &\sim \N{0}{1}, \label{eq:ds-imeas}\\
T^{\text{m}}_k &= T_k + \sigma_T\,\nu^T_k, & \nu^T_k &\sim \N{0}{1}, \label{eq:ds-tmeas}
\end{align}
where $v_k, i_k, T_k$ are the true terminal voltage, current and cell temperature;
$\mathcal{Q}_q(x) = q\,\mathrm{round}(x/q)$ is a mid-tread quantiser of step $q$;
$\kappa$ is the current-sensor gain (scale-factor) error; $b_i$ its offset;
$o_k$ an optional voltage outlier; and $d_k$ an optional stuck-at (dropout) fault. The nominal
values are collected in Table~\ref{tab:ds-sensors}.

\begin{table}[H]\centering\small
\caption{Nominal sensor model. These are the defaults of \code{struct Scenario}; the twelve
scenarios of \S\ref{sec:ds-scenarios} perturb a subset of them.}
\label{tab:ds-sensors}
\begin{tabular}{llll}\toprule
symbol & field & nominal value & meaning \\ \midrule
$\sigma_v$ & \code{sigma\_v} & \SI{2}{\milli\volt} & voltage-sensor noise (std) \\
$\sigma_i$ & \code{sigma\_i} & \SI{50}{\milli\ampere} & current-sensor noise (std) \\
$\sigma_T$ & \code{sigma\_T} & \SI{0.2}{\kelvin} & temperature-sensor noise (std) \\
$q_v$ & \code{v\_quant} & \SI{1}{\milli\volt} & voltage ADC quantisation step \\
$q_i$ & \code{i\_quant} & \SI{10}{\milli\ampere} & current ADC quantisation step \\
$\kappa$ & \code{current\_gain} & \num{1.005} & current-sensor scale-factor error ($+0.5\,\%$) \\
$b_i$ & \code{current\_bias} & \SI{20}{\milli\ampere} & current-sensor offset ($0.4\,\%$ of 1\,C) \\
$T_{\text{amb}}$, $T_0$ & \code{T\_amb}, \code{T0} & \SI{298.15}{\kelvin} & ambient and initial cell temperature \\
-- & \code{outlier\_prob} & 0 & probability per sample of a voltage outlier \\
-- & \code{v\_dropout\_start} & $-1$ (disabled) & start time of a stuck-at voltage fault \\
\bottomrule\end{tabular}
\end{table}

\subsection{Why the nominal case already has a biased current sensor}

The most consequential modelling decision in this chapter is that the \emph{nominal} scenario is not
a perfect current sensor. A $+0.5\,\%$ gain error and a \SI{20}{\milli\ampere} offset are
representative of a good, temperature-compensated, factory-calibrated Hall-effect or shunt sensor in
an airborne installation; they are not a fault. Making them part of the nominal case forces every
estimator to face, from the first row of every results table, the error mechanism that dominates
industrial SOC estimation. The reason is the structure of the Coulomb-counting integral: over a
mission of duration $T$ the error contributed by a gain error $\kappa-1$ and an offset $b_i$ is
\begin{equation}
\Delta \soc \;=\; -\frac{1}{3600\,Q}\left[(\kappa-1)\!\int_0^T\! i(t)\,\mathrm{d}t \;+\; b_i T\right]
 \;=\; -\frac{(\kappa-1)\,A_h + b_i T/3600}{Q},
\label{eq:ds-ccbias}
\end{equation}
with $A_h$ the charge actually removed in \si{\ampere\hour} and $Q$ the capacity. Note the
asymmetry: the gain error scales with \emph{throughput}, the offset with \emph{time}. For the eVTOL
mission ($A_h = \SI{3.763}{\ampere\hour}$, $T=\SI{2010}{\second}$, $Q=\SI{5}{\ampere\hour}$) the
gain term is \SI{0.376}{\percent} and the offset term \SI{0.223}{\percent}, a total of
\SI{0.60}{\percent} of SOC; the same offset would contribute \SI{0.32}{\percent} on the
\SI{48}{\minute} fixed-wing mission and \SI{0.13}{\percent} on the \SI{19}{\minute} hover-hold,
while the gain term barely changes. Measured over the nominal mission the difference between
integrated measured and integrated true current is \SI{0.0300}{\ampere\hour} --- exactly the
\SI{0.60}{\percent} predicted --- while the observed final Coulomb-counting error is
$-\SI{0.44}{\percent}$; the missing \SI{0.16}{\percent} is the in-flight capacity fade, which has
the opposite sign (Chapter~\ref{ch:coulombcounting} works the decomposition out in full).

This follows Ng et~al.~\cite{ng2009coulomb}, who showed that Coulomb counting is limited in practice
by sensor calibration and capacity knowledge rather than by the integration, and Lin
et~al.~\cite{lin2018sensorbias}, who showed that a current-sensor bias is \emph{unobservable} from a
single voltage measurement unless it is explicitly added to the state. A benchmark whose nominal
case had a perfect current sensor would rank Coulomb counting far too highly and give no credit to
the unknown-input observers (Chapter~\ref{ch:uio}), disturbance observers (Chapter~\ref{ch:dob}) and
augmented-state filters (Chapter~\ref{ch:jointekf}) that exist to recover it.

\subsection{Quantisation}

Quantisation is applied last, after the noise, so it is the ADC of a real BMS rather than an
independent noise source. With a uniform quantiser of step $q$ and an input whose noise is
comparable to or larger than $q$, the error is well modelled as uniform on $[-q/2,q/2]$ with
variance $q^2/12$~\cite{widrow2008quantization}: $q_v/\sqrt{12} = \SI{0.29}{\milli\volt}$ against
$\sigma_v = \SI{2}{\milli\volt}$, and $q_i/\sqrt{12} = \SI{2.9}{\milli\ampere}$ against
$\sigma_i = \SI{50}{\milli\ampere}$. Neither dominates, which is the intent: quantisation is
present, visible in the residual traces as a staircase at low current, and never the limiting error.
It does make the measured signals \emph{non-Gaussian} in a way that matters for the Student-$t$ and
maximum-correntropy filters of Part~V, whose scale estimates would otherwise be driven purely by
$\sigma_v$.

\section{The twelve cell-level scenarios}
\label{sec:ds-scenarios}

A scenario is a set of deviations from the nominal sensor model, the nominal plant, or the
parameter set handed to the estimator. The twelve scenarios returned by
\code{standard\_scenarios()} are listed in Table~\ref{tab:ds-scenarios}; every estimator is run on
every one of them, on both plants, so the results tables in Parts~II--X are directly comparable
row by row. Table~\ref{tab:ds-scenario-map} is the reader's map: it states what each scenario is
designed to expose and which estimator families are expected to react to it.

\begin{table}[H]\centering\small
\caption{The twelve cell-level scenarios (\code{standard\_scenarios()},
\code{battery/datasets.hpp}). Only the fields that differ from Table~\ref{tab:ds-sensors} are
listed; everything else keeps its nominal value. ``truth'' marks a change to the plant, ``est.''
a change to what the estimator is told.}
\label{tab:ds-scenarios}
\begin{tabular}{@{}llp{0.52\textwidth}@{}}\toprule
\# & scenario & deviation from nominal \\ \midrule
1 & \code{nominal} & --- (Table~\ref{tab:ds-sensors} as it stands: $\kappa=1.005$, $b_i=\SI{20}{\milli\ampere}$) \\
2 & \code{init\_error} & est.: $\soc_0^{\text{est}} = \soc_0^{\text{true}} - 0.35$ (initial SOC handed to the estimator is $0.60$ instead of $0.95$) \\
3 & \code{current\_bias} & $b_i = \SI{0.25}{\ampere}$ (5\,\% of 1\,C), $\kappa = 1.02$ ($+2\,\%$) \\
4 & \code{noise\_step} & $\sigma_v$ steps from \SI{2}{\milli\volt} to \SI{20}{\milli\volt} at $t=\SI{600}{\second}$ \\
5 & \code{outliers} & voltage outliers of $\pm\SI{50}{\milli\volt}$ (random sign) with probability \num{0.05} per sample \\
6 & \code{cold} & $T_{\text{amb}} = T_0 = \SI{263.15}{\kelvin}$ ($-\SI{10}{\celsius}$) \\
7 & \code{hot} & $T_{\text{amb}} = T_0 = \SI{318.15}{\kelvin}$ ($+\SI{45}{\celsius}$) \\
8 & \code{aged\_cell} & truth: \SI{15}{\percent} capacity loss pre-applied ($Q=\SI{4.25}{\ampere\hour}$, $R_0$ raised by $2.5\times$ the loss fraction, i.e.\ $+\SI{37.5}{\percent}$); est.\ still told \SI{5}{\ampere\hour} \\
9 & \code{param\_mismatch} & est.: $R_0 \times 0.7$, $R_1 \times 1.3$, $\tau_2 \times 1.5$ \\
10 & \code{sensor\_dropout} & voltage channel stuck at its last value for \SI{15}{\second} from $t=\SI{200}{\second}$ \\
11 & \code{preisach} & truth: the one-state hysteresis of the plant is replaced by a 32-level Preisach operator with $M_p=\SI{15}{\milli\volt}$, oriented as in a real cell (the \emph{charging} branch lies above the \emph{discharging} branch at the same SOC); est.\ still uses the one-state model with $M=\SI{5}{\milli\volt}$ \\
12 & \code{combined} & est.: $\soc_0$ low by \num{0.25}; $b_i=\SI{0.15}{\ampere}$, $\kappa=1.02$; truth: \SI{10}{\percent} capacity loss; $T_{\text{amb}}=T_0=\SI{273.15}{\kelvin}$ ($\SI{0}{\celsius}$) \\
\bottomrule\end{tabular}
\end{table}

\begin{table}[H]\centering\small
\caption{What each scenario is for. ``Expected responders'' names the estimator families whose
defining mechanism addresses that failure mode; a family that does \emph{not} improve on its own
scenario is a negative result worth reporting, and several such cases appear in Parts~IV--X.}
\label{tab:ds-scenario-map}
\begin{tabular}{@{}lp{0.40\textwidth}p{0.36\textwidth}@{}}\toprule
scenario & failure mode exposed & expected responders \\ \midrule
\code{nominal} & noise floor; steady-state bias from the calibrated sensor & all; separates numerically equivalent implementations \\
\code{init\_error} & transient behaviour from a poor prior; over-confident linearisation & anything with a covariance; iterated (IEKF, IPLF) and sigma-point filters; particle filters \\
\code{current\_bias} & unobservable input disturbance integrated by the model & unknown-input and disturbance observers, PI observers, augmented-state (joint/dual) filters, sliding-mode observers \\
\code{noise\_step} & $\mR$ wrong by $100\times$ half-way through the flight & innovation-adaptive KFs (Sage--Husa, IAE, fuzzy, variational Bayes), IMM/MMAE, $H_\infty$ \\
\code{outliers} & heavy-tailed measurement noise & Huber, Student-$t$, maximum-correntropy, SVSF, gated filters, particle filters \\
\code{cold} & $R_0$ up $2.1\times$, $\tau_2$ up $2.5\times$; large, slow overpotential & temperature-scheduled models; adaptive and joint-parameter filters \\
\code{hot} & impedance down $\approx 40\,\%$; fastest dynamics & same as \code{cold}, opposite sign; also a numerical-conditioning check \\
\code{aged\_cell} & capacity and resistance both wrong by a lot & joint/dual state--parameter estimation (Part~VII), AWTLS, RLS \\
\code{param\_mismatch} & structural model error with the right capacity & adaptive, robust ($H_\infty$, SVSF) and moving-horizon estimators \\
\code{sensor\_dropout} & measurement frozen: innovations become deterministic & residual-gating filters, interval observers, anything that degrades gracefully to open loop \\
\code{preisach} & unmodelled path-dependent hysteresis & multiple-model (IMM/MMAE/GSF), particle filters, hysteresis-augmented models \\
\code{combined} & four faults at once (prior, bias, ageing, cold) & the scenario that separates the field: see Chapter~\ref{ch:ekf} for why the plain EKF fails it \\
\bottomrule\end{tabular}
\end{table}

Three scenarios deserve an explanatory note.

\emph{\code{init\_error} is a $-35\,\%$ error, not $-5\,\%$.} An avionics BMS re-powered in flight,
or one that has lost its non-volatile SOC store, starts with essentially no information. The
benchmark uses $\soc_0^{\text{est}} = 0.60$ against a truth of \num{0.95}, placing the estimate on
the far side of the OCV inflection at $z\approx0.65$. Since the prior standard deviation handed to
every estimator is $\sigma_{\soc_0}=0.10$ regardless of scenario, the true error is $3.5\sigma$ and
the filter is by construction over-confident --- which is what distinguishes a filter whose
linearisation is valid over its own prior covariance from one whose is not.

\emph{\code{aged\_cell} changes the plant only.} The \SI{15}{\percent} capacity loss is applied to
the truth through \code{set\_capacity\_loss}, which also raises $R_0$ by
$2.5\times0.15 = \SI{37.5}{\percent}$ through the coupling used by the ageing law
(Chapter~\ref{ch:ecm}), while the estimator is still told $Q=\SI{5}{\ampere\hour}$ and
$R_0=\SI{12}{\milli\ohm}$. This is the honest version of the ageing problem: capacity and resistance
drift \emph{together}, so a filter that estimates only one carries the other as a bias.

\emph{\code{combined} is not the sum of its parts.} At $\SI{0}{\celsius}$, $R_0$ rises to
\SI{25.1}{\milli\ohm} and $\tau_2$ to \SI{302}{\second}; with a \SI{10}{\percent} capacity loss the
effective $R_0$ is another \SI{25}{\percent} higher, so at the \SI{22.5}{\ampere} hover the ohmic
drop alone exceeds \SI{0.7}{\volt}. The filter must simultaneously reject a $-0.25$ prior error, an
unobservable \SI{0.15}{\ampere} bias and a $\sim\SI{0.2}{\volt}$ modelling residual. Roughly half
the library fails it, and by the same mechanism in almost every case: premature covariance
collapse.

\section{What the estimator is told, and what it is not}
\label{sec:ds-truth}

At \code{reset()} every estimator receives one \code{EstimatorConfig}
(\code{battery/ecm\_model.hpp}) and nothing else.

\begin{table}[H]\centering\small
\caption{\code{EstimatorConfig}: the complete interface between the benchmark and an estimator.}
\label{tab:ds-cfg}
\begin{tabular}{@{}lll@{}}\toprule
field & nominal value & note \\ \midrule
\code{params} & \code{CellParams::nominal()} & the full ECM parameter set, possibly scaled by the scenario \\
\code{dt} & \SI{0.1}{\second} & sample time \\
\code{z0} & $\soc_0^{\text{true}} + \Delta$, clamped to $[0.02,1]$ & initial SOC guess, $\Delta$ from the scenario \\
\code{sigma\_z0} & \num{0.1} & assumed initial SOC standard deviation --- \emph{always} \num{0.1} \\
\code{sigma\_v} & \SI{2}{\milli\volt} & assumed voltage-noise std (the scenario's \emph{true} value) \\
\code{sigma\_i} & \SI{50}{\milli\ampere} & assumed current-noise std \\
\code{sigma\_T} & \SI{0.2}{\kelvin} & assumed temperature-noise std \\
\code{T0} & \SI{298.15}{\kelvin} & initial temperature \\
\code{seed} & benchmark seed & for stochastic estimators (particle filters, ensemble filters) \\
\bottomrule\end{tabular}
\end{table}

The estimator is therefore told the sensor noise standard deviations \emph{as they are at the start
of the run}, the sample time, and a parameter set. It is \textbf{not} told: the true SOC at any
time; the true capacity or its fade; the current-sensor gain $\kappa$ or offset $b_i$ (in any
scenario, \code{current\_bias} included); that the voltage noise steps up at \SI{600}{\second} in
\code{noise\_step}; that outliers are present; that the voltage channel will freeze; that the plant
uses a Preisach operator; or that the parameters it was handed are wrong. In
\code{param\_mismatch} it receives the \emph{scaled} parameters with no flag; in \code{aged\_cell},
the healthy ones.

Two deliberate concessions are worth flagging. First, \code{sigma\_v} is the true pre-step value
in \code{noise\_step}, which makes that scenario a clean test of \emph{adaptation} rather than of
mistuning. Second, the temperature is measured, not estimated: $u=[i,T]^\T$ carries the measured
cell temperature into every model evaluation, so temperature-dependent parameters are scheduled on
a noisy but unbiased signal. A production BMS measures surface temperature and infers core
temperature; modelling that gap would be a second, separate benchmark.

\section{The pack dataset}
\label{sec:ds-pack}

A flight pack is a series--parallel assembly, and its SOC is not a scalar. The pack dataset
(\code{pack/pack\_dataset.hpp}) models a module of $N_c = 12$ cells in series that share one
current and are each measured by their own voltage and temperature channel.

\subsection{Cell-to-cell spread}

Each cell $c$ is an independent \code{EcmPlant} whose parameters are drawn once from a manufacturing
spread:
\begin{align}
Q^{(c)} &= Q_{\text{nom}}\bigl(1 + s_Q\,\xi^{(c)}_Q\bigr), &
R^{(c)}_j &= R_j\bigl(1 + s_R\,\xi^{(c)}_R\bigr)\ \ (j=0,1,2), \label{eq:ds-pack-par}\\
\soc^{(c)}_0 &= \soc_0 + s_z\,\xi^{(c)}_z, &
T^{(c)}_{\text{amb}} &= T_{\text{amb}} + s_T\,\xi^{(c)}_T, \label{eq:ds-pack-init}
\end{align}
with all $\xi\sim\N{0}{1}$ independent and nominal spreads $s_Q = \SI{2}{\percent}$,
$s_R = \SI{8}{\percent}$ (the three resistances of a cell share \emph{one} draw, because they are
physically correlated), $s_z = \num{0.015}$ and $s_T = \SI{1.5}{\kelvin}$ (the thermal gradient
across the module). One current is applied to the string; per-cell voltages are measured with
$\sigma_v=\SI{2}{\milli\volt}$ and \SI{1}{\milli\volt} quantisation, and the shared current with the
nominal model of Table~\ref{tab:ds-sensors}. At seed~1 the draw gives capacities between
\SIlist{4.79;5.12}{\ampere\hour} and initial SOCs between \numlist{0.933;0.970}; after one eVTOL
mission the true cell-SOC spread has grown from \num{3.7} to \num{6.3} percentage points, purely
from the capacity spread, because the same charge removed from a smaller cell is a larger fraction
of it.

\subsection{Pack-level definitions of SOC}

A BMS must report three different numbers, all of which are called ``pack SOC''
(Plett~\cite{plett2016bms2}, Vol.~II Ch.~4):
\begin{align}
\soc^{\text{pack}}_{\min} &= \min_c \soc^{(c)}, &&\text{usable discharge --- the weakest cell hits the lower limit first;}\label{eq:ds-socmin}\\
\soc^{\text{pack}}_{\max} &= \max_c \soc^{(c)}, &&\text{usable charge --- the strongest cell hits the upper limit first;}\label{eq:ds-socmax}\\
\soc^{\text{pack}}_{\text{mean}} &= \tfrac{1}{N_c}\textstyle\sum_c \soc^{(c)}, &&\text{energy accounting.}\label{eq:ds-socmean}
\end{align}
These are estimated quantities in their own right and are \emph{not} interchangeable: a method
excellent at $\soc^{\text{pack}}_{\text{mean}}$ can be poor at $\soc^{\text{pack}}_{\min}$, because
averaging cancels the per-cell errors that the minimum selects for. All three are scored separately
(Chapter~\ref{ch:metrics}).

\subsection{The six pack scenarios}
\label{sec:ds-pack-scen}

\begin{table}[H]\centering\small
\caption{Pack scenarios (\code{pack\_scenarios()}). Nominal spreads: $s_Q=\SI{2}{\percent}$,
$s_R=\SI{8}{\percent}$, $s_z=\num{0.015}$, $s_T=\SI{1.5}{\kelvin}$.}
\label{tab:ds-pack-scen}
\begin{tabular}{@{}llp{0.46\textwidth}@{}}\toprule
\# & scenario & deviation and purpose \\ \midrule
1 & \code{nominal} & --- : baseline per-cell accuracy and the cost of $12\times$ everything \\
2 & \code{init\_error} & every cell's estimator starts \num{0.35} low: does the pack method converge as one, or cell by cell? \\
3 & \code{current\_bias} & $b_i = \SI{0.25}{\ampere}$, $\kappa=1.02$ on the \emph{shared} current sensor: a common-mode disturbance that a consensus or bar--delta method can in principle separate from per-cell error \\
4 & \code{weak\_cell} & cell~5 pre-aged by \SI{12}{\percent} capacity with its $R_0, R_1$ scaled by \num{1.3}: the single outlier that determines $\soc^{\text{pack}}_{\min}$ \\
5 & \code{noise\_step} & $\sigma_v$ steps to \SI{20}{\milli\volt} at \SI{600}{\second} on all channels \\
6 & \code{large\_imbalance} & $s_z = \num{0.05}$, $s_Q = \SI{5}{\percent}$: an unbalanced or partially repaired module \\
\bottomrule\end{tabular}
\end{table}

\code{weak\_cell} is the pack analogue of \code{combined}: at seed~1 cell~5 starts from
\SI{4.21}{\ampere\hour} against a module mean of \SI{4.92}{\ampere\hour}, and by the end of the
mission its true SOC is \num{0.066} against \num{0.215} for the strongest cell --- a spread of
\num{15} percentage points. Every cell is handed the \emph{same} nominal configuration, so a method
that assumes a homogeneous module reports the mean and misses the cell that limits the mission.

\estkitfig{dataset_pack.png}{The twelve-cell module on the eVTOL mission (nominal spread, seed~1):
true per-cell SOC, with the weakest and strongest cells highlighted. The fan opens monotonically
because the capacity spread, not the initial imbalance, dominates after the first
minutes.}{fig:ds-pack}

\section{The inertial flight dataset}
\label{sec:ds-imu}

The attitude part of the benchmark exists because the estimator families in this report were
invented for aerospace navigation, and a battery-only benchmark cannot show what a multiplicative or
invariant formulation buys. The dataset (\code{attitude/imu\_dataset.hpp}) is a \SI{900}{\second}
flight of a small fixed-wing aircraft at \SI{100}{\hertz} ($n=\num{90000}$ samples).

\subsection{Frames and conventions}

The navigation frame is NED (north, east, down) and the body frame FRD (forward, right, down), the
standard aerospace pair~\cite{groves2013,titterton2004}; the unit quaternion $q$ is Hamilton,
scalar-first, and maps \emph{body to navigation}, $v_{\text{nav}} = R(q)\,v_{\text{body}}$; Euler
angles are the ZYX (yaw--pitch--roll) aerospace sequence~\cite{kuipers1999}; and gravity in the
navigation frame is $g_{\text{nav}} = [0,0,+g]^\T$ with
$g = \SI{9.80665}{\metre\per\second\squared}$. The three sensors are modelled as
\begin{align}
\omega^{\text{m}}_k &= \omega_k + b^g_k + n^g_k, &
 b^g_k &= b^g_{k-1} + w^g_k, \label{eq:ds-gyro}\\
f^{\text{m}}_k &= R(q_k)^\T\bigl(a^{\text{nav}}_k - g_{\text{nav}}\bigr) + b^a + n^a_k, \label{eq:ds-accel}\\
m^{\text{m}}_k &= R(q_k)^\T m_{\text{nav}} + b^m + n^m_k, \label{eq:ds-mag}
\end{align}
where $\omega_k$ is the true body rate, $b^g$ a random-walk gyro bias, $a^{\text{nav}}$ the inertial
acceleration in the navigation frame, and $b^a$, $b^m$ constant accelerometer and magnetometer
biases. Equation~\eqref{eq:ds-accel} is the definition of \emph{specific force}: at rest the
accelerometer reads $R^\T[0,0,-g]^\T$, i.e.\ $+g$ along body $-z$ --- a sign convention that is a
recurring source of error between implementations, revisited in Chapter~\ref{ch:metrics}. The
reference field is $|m_{\text{nav}}| = \SI{48}{\micro\tesla}$ at inclination
$I=\SI{65}{\degree}$, representative of central Europe~\cite{chulliat2020wmm}, giving
$m_{\text{nav}} = [\,48\cos I,\ 0,\ 48\sin I\,]^\T = [\,20.29,\ 0,\ 43.50\,]^\T\,\si{\micro\tesla}$
(declination is taken as zero, so magnetic and true north coincide). A residual hard-iron offset
$b^m = [\,0.30,\ -0.15,\ 0.25\,]^\T\,\si{\micro\tesla}$, of magnitude
$\norm{b^m} = \SI{0.42}{\micro\tesla}$ --- about \SI{0.8}{\percent} of the field --- represents a
compass that has been calibrated but not perfectly. That figure is deliberate: a magnetometer whose
hard-iron error had been left at the several-microtesla level of an uncalibrated installation would
make the \code{mag\_disturbance} scenario indistinguishable from the nominal one, and would swamp
the yaw performance of every filter in Part~XII with a constant offset rather than testing their
response to a transient.

\subsection{The scripted mission}

The truth attitude is generated from a script of Euler angles and airspeed, not from a dynamic
model:

\begin{table}[H]\centering\small
\caption{Scripted flight phases of the inertial dataset. Turbulence (\S\ref{sec:ds-imu-turb}) is
superimposed on all three angles, and the whole script is then low-pass filtered.}
\label{tab:ds-imu-script}
\begin{tabular}{@{}llrrrl@{}}\toprule
$t$ [\si{\second}] & phase & roll [\si{\degree}] & pitch [\si{\degree}] & yaw rate [\si{\degree\per\second}] & airspeed [\si{\metre\per\second}] \\ \midrule
0--60    & take-off roll   & 0    & 0   & 0    & $20 \to 44$ (ramp) \\
60--180  & climb           & 0    & 10  & 0    & 44 \\
180--240 & right turn      & $+25$& 4   & $+5.2$ & 45 \\
240--480 & cruise          & 0    & 2   & 0    & 45 \\
480--540 & left turn       & $-25$& 3   & $-5.2$ & 45 \\
540--780 & descent         & 0    & $-4$& 0    & 40 \\
780--840 & base turn       & $+20$& 1   & $+4.0$ & 45 \\
840--900 & final approach  & 0    & $-3$& 0    & 35 \\
\bottomrule\end{tabular}
\end{table}

\subsection{From a script to sensor readings}
\label{sec:ds-imu-turb}

Four steps turn Table~\ref{tab:ds-imu-script} into \eqref{eq:ds-gyro}--\eqref{eq:ds-mag}.

\emph{(i) Turbulence.} Deterministic sinusoids are added to each angle: $\pm\SI{3}{\degree}$ at
\SI{0.5}{\hertz} plus $\pm\SI{1.5}{\degree}$ at \SI{0.13}{\hertz} on roll, $\pm\SI{1.5}{\degree}$ at
\SI{0.3}{\hertz} plus $\pm\SI{0.7}{\degree}$ at \SI{0.07}{\hertz} on pitch, $\pm\SI{1.0}{\degree}$ at
\SI{0.05}{\hertz} on yaw, all scaled by a factor the \code{high\_dynamics} scenario raises to~3.
Fixed sinusoids rather than noise keep the truth exactly reproducible and put the disturbance energy
at known frequencies, so a filter's bandwidth can be read off its error spectrum.

\emph{(ii) Low-pass.} The scripted angles pass through a first-order lag,
$x_k = x_{k-1} + (\dt/\tau)(x^{\text{script}}_k - x_{k-1})$ with $\tau=\SI{2}{\second}$, so the
phase changes are differentiable and the body rates finite; yaw is unwrapped first.

\emph{(iii) Body rates.} The Euler rates $(\dot\phi,\dot\theta,\dot\psi)$ come from central
differences of the filtered angles and are converted to body rates by the inverse ZYX Euler-rate
kinematics
\begin{equation}
\omega \;=\;
\begin{bmatrix} \dot\phi - \dot\psi\sin\theta \\
 \dot\theta\cos\phi + \dot\psi\cos\theta\sin\phi \\
 \dot\psi\cos\theta\cos\phi - \dot\theta\sin\phi \end{bmatrix},
\label{eq:ds-eulerrate}
\end{equation}
which is exact for the ZYX sequence and is the standard result~\cite{stevens2015aircraft}.

\emph{(iv) Specific force.} The navigation-frame flight-path velocity is built from the airspeed
$V$, the heading $\psi$ and a flight-path angle $\gamma = \theta - \SI{3}{\degree}$ (a constant
angle of attack):
\begin{equation}
v_{\text{nav}} = V\bigl[\cos\psi\cos\gamma,\ \sin\psi\cos\gamma,\ -\sin\gamma\bigr]^\T,
\qquad a^{\text{nav}}_k = \frac{v_{\text{nav},k} - v_{\text{nav},k-1}}{\dt},
\label{eq:ds-flightpath}
\end{equation}
and \eqref{eq:ds-accel} is evaluated with that acceleration. Building the acceleration from the
kinematics rather than prescribing it makes the dataset \emph{self-consistent}: the turns carry
exactly the centripetal acceleration their turn rate and airspeed imply.

\subsection{Sensor noise and the six scenarios}

Nominal sensor parameters (all per sample at \SI{100}{\hertz}) are: gyro noise
$\sigma_g = \SI{0.1}{\degree\per\second}$; initial gyro bias magnitude
\SI{0.5}{\degree\per\second} per axis, applied as
$b^g_0 = [\,+1,\,-0.6,\,+0.4\,]\times\SI{0.5}{\degree\per\second}$; gyro-bias random walk
\SI{0.002}{\degree\per\second} per $\sqrt{\text{sample}}$; accelerometer noise
$\sigma_a = \SI{0.05}{\metre\per\second\squared}$ and bias
$b^a = [\,+1,\,-1,\,+0.5\,]\times\SI{0.02}{\metre\per\second\squared}$; magnetometer noise
$\sigma_m = \SI{0.5}{\micro\tesla}$.

\begin{table}[H]\centering\small
\caption{The six attitude scenarios (\code{imu\_scenarios()}).}
\label{tab:ds-imu-scen}
\begin{tabular}{@{}llp{0.50\textwidth}@{}}\toprule
\# & scenario & deviation and purpose \\ \midrule
1 & \code{nominal} & --- : bias observability and steady-state accuracy \\
2 & \code{init\_error} & the initial estimate is rotated by \SI{30}{\degree} about \emph{each} body axis \\
3 & \code{gyro\_bias} & initial bias raised to \SI{3}{\degree\per\second} per axis ($6\times$) \\
4 & \code{mag\_disturbance} & \SI{25}{\micro\tesla} added along navigation east for \SI{20}{\second} from $t=\SI{300}{\second}$ (over half the field) \\
5 & \code{vibration} & accelerometer noise raised by \SI{1.0}{\metre\per\second\squared} (rotor or propeller vibration) \\
6 & \code{high\_dynamics} & turbulence amplitude $\times 3$ \\
\bottomrule\end{tabular}
\end{table}

\paragraph{A \SI{30}{\degree} per-axis error is a \SI{46.6}{\degree} error.} The initial estimate is
formed as $\hat q_0 = q_0 \otimes q_{\mathrm{euler}}(30^\circ,30^\circ,30^\circ)$. Rotations do
not add: the scalar part of $q_{\mathrm{euler}}(30^\circ,30^\circ,30^\circ)$ is
$\cos^3\!15^\circ + \sin^3\!15^\circ = 0.91856$, so the total angular error is
$2\arccos(0.91856) = \SI{46.62}{\degree}$. Quoting the per-axis figure would understate the
difficulty by more than a third, and the angular distance \eqref{eq:ds-imu-angdist} is what the
benchmark actually scores.

\paragraph{Coordinated turns defeat gravity referencing.} This is the most instructive feature of
the dataset. In a steady coordinated turn at bank angle $\phi$ the horizontal (centripetal)
component of the specific force is $g\tan\phi$; for $\phi=\SI{25}{\degree}$ that is
$\SI{4.57}{\metre\per\second\squared} = \SI{0.47}{}g$, sustained for the whole \SI{60}{\second} of
the turn. The vector sum of gravity and that horizontal acceleration lies \emph{along the body $z$
axis} --- exactly what the accelerometer reads in level flight. A gravity-referenced levelling
algorithm (the accelerometer correction of a complementary, Madgwick or Mahony filter, or the
accelerometer update of a multiplicative EKF that models $f = -R^\T g$) therefore reads a
coordinated turn as \emph{zero roll} and drives the estimate towards it with whatever gain it has.
In the dataset the scripted turn rate of \SI{5.2}{\degree\per\second} at
\SI{45}{\metre\per\second} is close to, but not exactly, the \SI{5.82}{\degree\per\second} that
would make the \SI{25}{\degree} bank truly coordinated: the measured horizontal specific force
during the right turn averages \SI{4.08}{\metre\per\second\squared} ($\SI{0.42}{}g$) and the total
body-frame specific force reaches $\SI{1.17}{}g$. The mismatch is small enough that the trap is
fully sprung, which is why every gravity-referenced filter in Part~XII shows a roll error that grows
through the turns and decays afterwards, why estimators that gate the accelerometer when
$\bigl|\,|f^{\text{m}}| - g\,\bigr|$ is large do measurably better, and why the cross-validation of
Chapter~\ref{ch:metrics} reports attitude RMS errors of \SIrange{17}{23}{\degree} for ungated
Madgwick and Mahony --- those are not implementation defects, they are the turns.

\estkitfig{dataset_imu.png}{The synthetic flight used by the attitude benchmark (nominal, seed~1):
truth Euler angles, gyroscope, accelerometer (specific force) and magnetometer. The two
\SI{25}{\degree} turns at \SIlist{3;8}{\minute} are visible in every channel; the magnetometer
traces show the field rotating through the body frame as heading changes by
\SI{312}{\degree}.}{fig:ds-imu}

\section{The multi-flight state-of-health protocol}
\label{sec:ds-soh}

Capacity fade is a phenomenon of hundreds of cycles, so state-of-health (SOH) estimation cannot be
benchmarked on a single mission. \code{bench\_soh} runs a fleet-operations protocol in which the
\emph{estimators run continuously} across the whole study while the plant ages underneath them. One
cycle --- nominally one day of operations --- has four phases:
\begin{enumerate}[leftmargin=*,topsep=2pt,itemsep=1pt]
\item \textbf{Flight}: the chosen mission profile (eVTOL by default), \SI{2010}{\second};
\item \textbf{Rest}: \SI{1200}{\second} (\SI{20}{\minute}) at zero current --- turnaround;
\item \textbf{CC--CV recharge}: 1\,C constant current to \SI{4.20}{\volt}, then a
voltage-limited taper, terminating at $C/20$ or $\soc=0.999$, capped at \SI{5400}{\second};
\item \textbf{Rest}: a further \SI{1200}{\second}.
\end{enumerate}
Between cycles the plant is advanced by one day of calendar ageing (\code{age\_calendar(1)}), which
applies the $\sqrt{t}$ storage term and relaxes the RC and thermal states to ambient. The study runs
\textbf{120 flights} --- 120 simulated days and roughly \num{450}\,\si{\ampere\hour} of throughput.

Cycle ageing at the true rate would move the capacity by only a per cent over 120 cycles, below the
resolution of most estimators, so the protocol applies an \textbf{ageing acceleration factor} $A=3$
(the \code{--accel} flag): at every step the plant's capacity-loss increment is multiplied by $A$
before being written back. This scales the \emph{rate} without changing its functional form
(Arrhenius $\times$ power law in throughput), so the shape a filter must track --- fast early fade,
slowing as $A_h^{z-1}$ with $z=0.55$ --- is preserved.

Only estimators reporting a finite capacity take part; \code{bench\_soh} discovers them at start-up
by resetting each one, stepping it twice and testing \code{capacity\_Ah()}. What is scored, per
flight, is:
\begin{itemize}[leftmargin=*,topsep=2pt,itemsep=1pt]
\item $\hat Q_f$, the capacity estimate averaged over the last \SI{10}{\percent} of the flight
phase (not over the rest or the charge, where the excitation is different);
\item $Q_f$, the true capacity, sampled at the quarter point of the cycle;
\item the per-flight SOC RMSE over the flight phase only;
\item the cumulative \si{\ampere\hour} throughput, as the independent variable for fade curves.
\end{itemize}
Capacity RMSE and MAE are computed over flights \numrange{11}{120}, discarding the first ten so that
the initial convergence of the parameter estimate does not dominate a metric meant to describe
tracking. The final capacity error and the mean flight SOC RMSE are reported separately, because an
SOH estimator that buys a good capacity number with SOC accuracy has solved the wrong problem.

\section{Reproducibility and summary}

Every dataset here is a pure function of its seed. The load, sensor, cell-spread and inertial
streams are separate \code{Rng} instances with different seed derivations
(Appendix~\ref{app:params}, Table~\ref{tab:ap-seeds}), so changing one never perturbs another; the
generator is xoshiro256$^{**}$~\cite{blackman2021xoshiro} with Marsaglia-polar
normals~\cite{marsaglia1964polar}, producing identical bit streams on every platform. Because the
seed drives the \emph{load} as well as the noise, the three seeds of the cell and pack benchmarks
are three different flights, not three noise realisations of one flight.

The data are built to make specific failure modes visible rather than to be realistic in the
aggregate. Four profiles span the load space from a \SI{25}{\minute} \SI{4}{\ampere} cruise to a
\SI{22.5}{\ampere} hover and a CC--CV recharge; a deliberately imperfect but non-faulty sensor chain
puts a \SI{0.6}{\percent} Coulomb-counting drift into the nominal case; twelve scenarios switch on
one pathology at a time and a thirteenth switches on four; a twelve-cell module adds the distinction
between mean, weakest and strongest cell; a \SI{15}{\minute} inertial flight adds the
coordinated-turn ambiguity that no amount of filtering can resolve from gravity alone; and a
120-flight protocol adds the slow drift that only joint state--parameter methods can follow.
Chapter~\ref{ch:metrics} defines what is measured on these data, and how the measurements were
checked against independent implementations.
